\documentclass[11pt]{article}
\newcommand{\ListaTitulo}{Lista 08 --- ANCOVA e Kruskal-Wallis}
% Preâmbulo compartilhado — MATD48, Listas de Exercícios 2026
% Usado por Lista{NN}.tex e Gabarito{NN}.tex via % Preâmbulo compartilhado — MATD48, Listas de Exercícios 2026
% Usado por Lista{NN}.tex e Gabarito{NN}.tex via % Preâmbulo compartilhado — MATD48, Listas de Exercícios 2026
% Usado por Lista{NN}.tex e Gabarito{NN}.tex via \input{preamble}

\usepackage[utf8]{inputenc}
\usepackage[T1]{fontenc}
\usepackage[brazil]{babel}
\usepackage{fullpage}
\usepackage{amsmath,amssymb}
\usepackage{enumitem}
\usepackage{xcolor}
\usepackage{fancyhdr}
\usepackage{titlesec}
\usepackage{listings}
\usepackage{hyperref}
\usepackage{booktabs}
\usepackage{graphicx}

\definecolor{matd48azul}{HTML}{0B4F6C}
\definecolor{matd48verde}{HTML}{2E8B57}

\titleformat{\section}{\normalfont\Large\bfseries\color{matd48azul}}{\thesection}{1em}{}
\titleformat{\subsection}{\normalfont\large\bfseries\color{matd48azul}}{\thesubsection}{1em}{}

\providecommand{\ListaTitulo}{Lista de Exercícios}

\setlength{\headheight}{14pt}
\pagestyle{fancy}
\fancyhf{}
\lhead{\small MATD48 --- Planejamento de Experimentos A}
\rhead{\small \ListaTitulo}
\cfoot{\thepage}
\renewcommand{\headrulewidth}{0.4pt}

\lstset{
  basicstyle=\ttfamily\small,
  breaklines=true,
  frame=single,
  columns=fullflexible,
  backgroundcolor=\color{gray!8},
  commentstyle=\color{matd48verde},
  keywordstyle=\color{matd48azul}\bfseries,
  language=R,
  extendedchars=true,
  literate=%
    {á}{{\'a}}1 {à}{{\`a}}1 {â}{{\^a}}1 {ã}{{\~a}}1
    {é}{{\'e}}1 {ê}{{\^e}}1 {í}{{\'i}}1 {ó}{{\'o}}1
    {ô}{{\^o}}1 {õ}{{\~o}}1 {ú}{{\'u}}1 {ç}{{\c c}}1
    {Á}{{\'A}}1 {À}{{\`A}}1 {Â}{{\^A}}1 {Ã}{{\~A}}1
    {É}{{\'E}}1 {Ê}{{\^E}}1 {Í}{{\'I}}1 {Ó}{{\'O}}1
    {Ô}{{\^O}}1 {Õ}{{\~O}}1 {Ú}{{\'U}}1 {Ç}{{\c C}}1
}

\hypersetup{colorlinks=true, linkcolor=matd48azul, urlcolor=matd48azul, citecolor=matd48azul}

\newcommand{\cabecalho}[2]{%
  \begin{center}
    {\Large\bfseries\color{matd48azul} #1}\\[0.3em]
    {\large #2}\\[0.3em]
    \rule{\linewidth}{0.6pt}
  \end{center}
  \vspace{0.5em}
}

\newenvironment{questao}[1]{\vspace{1em}\noindent\textbf{Questão #1.}\hspace{0.4em}}{\par}
\newenvironment{solucao}{\vspace{0.3em}\noindent\textit{Solução.}\hspace{0.4em}\color{black!85}}{\par\vspace{0.5em}}


\usepackage[utf8]{inputenc}
\usepackage[T1]{fontenc}
\usepackage[brazil]{babel}
\usepackage{fullpage}
\usepackage{amsmath,amssymb}
\usepackage{enumitem}
\usepackage{xcolor}
\usepackage{fancyhdr}
\usepackage{titlesec}
\usepackage{listings}
\usepackage{hyperref}
\usepackage{booktabs}
\usepackage{graphicx}

\definecolor{matd48azul}{HTML}{0B4F6C}
\definecolor{matd48verde}{HTML}{2E8B57}

\titleformat{\section}{\normalfont\Large\bfseries\color{matd48azul}}{\thesection}{1em}{}
\titleformat{\subsection}{\normalfont\large\bfseries\color{matd48azul}}{\thesubsection}{1em}{}

\providecommand{\ListaTitulo}{Lista de Exercícios}

\setlength{\headheight}{14pt}
\pagestyle{fancy}
\fancyhf{}
\lhead{\small MATD48 --- Planejamento de Experimentos A}
\rhead{\small \ListaTitulo}
\cfoot{\thepage}
\renewcommand{\headrulewidth}{0.4pt}

\lstset{
  basicstyle=\ttfamily\small,
  breaklines=true,
  frame=single,
  columns=fullflexible,
  backgroundcolor=\color{gray!8},
  commentstyle=\color{matd48verde},
  keywordstyle=\color{matd48azul}\bfseries,
  language=R,
  extendedchars=true,
  literate=%
    {á}{{\'a}}1 {à}{{\`a}}1 {â}{{\^a}}1 {ã}{{\~a}}1
    {é}{{\'e}}1 {ê}{{\^e}}1 {í}{{\'i}}1 {ó}{{\'o}}1
    {ô}{{\^o}}1 {õ}{{\~o}}1 {ú}{{\'u}}1 {ç}{{\c c}}1
    {Á}{{\'A}}1 {À}{{\`A}}1 {Â}{{\^A}}1 {Ã}{{\~A}}1
    {É}{{\'E}}1 {Ê}{{\^E}}1 {Í}{{\'I}}1 {Ó}{{\'O}}1
    {Ô}{{\^O}}1 {Õ}{{\~O}}1 {Ú}{{\'U}}1 {Ç}{{\c C}}1
}

\hypersetup{colorlinks=true, linkcolor=matd48azul, urlcolor=matd48azul, citecolor=matd48azul}

\newcommand{\cabecalho}[2]{%
  \begin{center}
    {\Large\bfseries\color{matd48azul} #1}\\[0.3em]
    {\large #2}\\[0.3em]
    \rule{\linewidth}{0.6pt}
  \end{center}
  \vspace{0.5em}
}

\newenvironment{questao}[1]{\vspace{1em}\noindent\textbf{Questão #1.}\hspace{0.4em}}{\par}
\newenvironment{solucao}{\vspace{0.3em}\noindent\textit{Solução.}\hspace{0.4em}\color{black!85}}{\par\vspace{0.5em}}


\usepackage[utf8]{inputenc}
\usepackage[T1]{fontenc}
\usepackage[brazil]{babel}
\usepackage{fullpage}
\usepackage{amsmath,amssymb}
\usepackage{enumitem}
\usepackage{xcolor}
\usepackage{fancyhdr}
\usepackage{titlesec}
\usepackage{listings}
\usepackage{hyperref}
\usepackage{booktabs}
\usepackage{graphicx}

\definecolor{matd48azul}{HTML}{0B4F6C}
\definecolor{matd48verde}{HTML}{2E8B57}

\titleformat{\section}{\normalfont\Large\bfseries\color{matd48azul}}{\thesection}{1em}{}
\titleformat{\subsection}{\normalfont\large\bfseries\color{matd48azul}}{\thesubsection}{1em}{}

\providecommand{\ListaTitulo}{Lista de Exercícios}

\setlength{\headheight}{14pt}
\pagestyle{fancy}
\fancyhf{}
\lhead{\small MATD48 --- Planejamento de Experimentos A}
\rhead{\small \ListaTitulo}
\cfoot{\thepage}
\renewcommand{\headrulewidth}{0.4pt}

\lstset{
  basicstyle=\ttfamily\small,
  breaklines=true,
  frame=single,
  columns=fullflexible,
  backgroundcolor=\color{gray!8},
  commentstyle=\color{matd48verde},
  keywordstyle=\color{matd48azul}\bfseries,
  language=R,
  extendedchars=true,
  literate=%
    {á}{{\'a}}1 {à}{{\`a}}1 {â}{{\^a}}1 {ã}{{\~a}}1
    {é}{{\'e}}1 {ê}{{\^e}}1 {í}{{\'i}}1 {ó}{{\'o}}1
    {ô}{{\^o}}1 {õ}{{\~o}}1 {ú}{{\'u}}1 {ç}{{\c c}}1
    {Á}{{\'A}}1 {À}{{\`A}}1 {Â}{{\^A}}1 {Ã}{{\~A}}1
    {É}{{\'E}}1 {Ê}{{\^E}}1 {Í}{{\'I}}1 {Ó}{{\'O}}1
    {Ô}{{\^O}}1 {Õ}{{\~O}}1 {Ú}{{\'U}}1 {Ç}{{\c C}}1
}

\hypersetup{colorlinks=true, linkcolor=matd48azul, urlcolor=matd48azul, citecolor=matd48azul}

\newcommand{\cabecalho}[2]{%
  \begin{center}
    {\Large\bfseries\color{matd48azul} #1}\\[0.3em]
    {\large #2}\\[0.3em]
    \rule{\linewidth}{0.6pt}
  \end{center}
  \vspace{0.5em}
}

\newenvironment{questao}[1]{\vspace{1em}\noindent\textbf{Questão #1.}\hspace{0.4em}}{\par}
\newenvironment{solucao}{\vspace{0.3em}\noindent\textit{Solução.}\hspace{0.4em}\color{black!85}}{\par\vspace{0.5em}}


\begin{document}

\cabecalho{Lista de Exercícios 08}{Análise de covariância e teste de Kruskal-Wallis (Capítulo 3, segunda metade / Aula 08)}

\noindent \textit{Entrega: até a próxima aula. Justifique todas as respostas e mostre os cálculos
intermediários --- nestas listas, a justificativa vale mais que a resposta final. O gabarito
completo está em \texttt{Gabarito08.pdf}.}

\begin{questao}{1} \textbf{(Psicologia --- médias ajustadas em ANCOVA)}

No experimento de dose de cafeína, a atenção basal (medida antes da dose) é usada como
covariável. O modelo ajustado forneceu inclinação comum $\hat\beta = 0{,}60$ e média geral da
covariável $\bar x_{\cdot\cdot} = 60$. Para dois dos cinco grupos de dose:

\begin{center}
\begin{tabular}{lccc}
\toprule
Dose & $\bar y_{i\cdot}$ (atenção pós-dose) & $\bar x_{i\cdot}$ (basal médio do grupo) \\
\midrule
0 mg   & 91{,}0 & 57 \\
200 mg & 108{,}0 & 64 \\
\bottomrule
\end{tabular}
\end{center}

\begin{enumerate}[label=(\alph*)]
  \item Calcule as médias ajustadas $\bar y_{i\cdot(\text{ajustada})} = \bar y_{i\cdot} -
    \hat\beta(\bar x_{i\cdot} - \bar x_{\cdot\cdot})$ para os dois grupos.
  \item Compare a diferença bruta $\bar y_{200} - \bar y_0$ com a diferença ajustada. O ajuste
    aumentou ou diminuiu a diferença estimada entre os grupos? Explique, em termos do
    desbalanceamento observado na covariável entre os grupos, por que isso aconteceu.
  \item Por que faz sentido reportar as médias \emph{ajustadas}, e não as médias brutas, como
    resultado principal de uma ANCOVA?
\end{enumerate}
\end{questao}

\begin{questao}{2} \textbf{(Ciência de dados --- verificando a homogeneidade das inclinações)}

Uma equipe testa quatro configurações de página quanto ao tempo de carregamento, usando o
número de imagens na página como covariável. Antes de rodar a ANCOVA, ela ajusta o modelo com a
interação covariável $\times$ configuração e obtém a seguinte tabela:

\begin{center}
\begin{tabular}{lccccc}
\toprule
Fonte & $gl$ & $SQ$ & $QM$ & $F$ & valor-$p$ \\
\midrule
n. imagens $\times$ configuração & 3 & 45{,}2 & 15{,}07 & 1{,}35 & 0{,}268 \\
\bottomrule
\end{tabular}
\end{center}

\begin{enumerate}[label=(\alph*)]
  \item O que exatamente está sendo testado por essa linha da tabela? Enuncie a hipótese nula em
    palavras e em termos do modelo $y_{ij} = \mu+\tau_i+\beta(x_{ij}-\bar x_{\cdot\cdot})+
    \varepsilon_{ij}$ (ou sua extensão com $\beta_i$ específico por grupo).
  \item Com base no valor-$p$ reportado, a equipe pode prosseguir com o modelo aditivo padrão de
    ANCOVA (inclinação única $\beta$)? Justifique.
  \item Suponha que o valor-$p$ tivesse sido $0{,}004$ em vez de $0{,}268$. O que a equipe deveria
    fazer de diferente, e por quê o ajuste por uma única inclinação comum deixaria de ser
    apropriado?
\end{enumerate}
\end{questao}

\begin{questao}{3} \textbf{(Agricultura --- Kruskal-Wallis, cálculo manual)}

Um agrônomo compara a produtividade (sacos/ha) de três variedades de feijão, com $n=3$ parcelas
por variedade, mas suspeita de forte assimetria nos dados e prefere não assumir normalidade:

\begin{center}
\begin{tabular}{lccc}
\toprule
Variedade A & Variedade B & Variedade C \\
\midrule
12 & 20 & 10 \\
18 & 25 & 14 \\
15 & 22 & 11 \\
\bottomrule
\end{tabular}
\end{center}

\begin{enumerate}[label=(\alph*)]
  \item Ordene as $N=9$ observações combinadas e atribua um posto (1 a 9) a cada uma. Calcule a
    soma dos postos $R_A$, $R_B$, $R_C$ por variedade, e verifique que $R_A+R_B+R_C = N(N+1)/2$.
  \item Calcule a estatística $H = \dfrac{12}{N(N+1)}\left(\dfrac{R_A^2}{3}+\dfrac{R_B^2}{3}+
    \dfrac{R_C^2}{3}\right) - 3(N+1)$.
  \item Compare $H$ com o valor crítico $\chi^2_{0{,}05;\,2} \approx 5{,}99$ e conclua.
\end{enumerate}
\end{questao}

\begin{questao}{4} \textbf{(Dedução --- não viesamento da média ajustada em ANCOVA)}

Considere o modelo de ANCOVA $y_{ij} = \mu + \tau_i + \beta(x_{ij}-\bar x_{\cdot\cdot}) +
\varepsilon_{ij}$, com $\varepsilon_{ij}$ de média zero e a covariável $x_{ij}$ tratada como
fixa (não aleatória, medida sem erro). Seja $\hat\beta$ um estimador não viesado de $\beta$
(isto é, $E[\hat\beta]=\beta$), obtido de forma independente da média amostral $\bar y_{i\cdot}$
do grupo $i$.

\begin{enumerate}[label=(\alph*)]
  \item Mostre que, sob o modelo, $E[\bar y_{i\cdot}] = \mu_i + \beta(\bar x_{i\cdot} - \bar
    x_{\cdot\cdot})$, em que $\mu_i = \mu+\tau_i$ é a média populacional (ajustada) do grupo $i$.
  \item Usando o item (a) e $E[\hat\beta]=\beta$, mostre que a média ajustada
    $\bar y_{i\cdot(\text{ajustada})} = \bar y_{i\cdot} - \hat\beta(\bar x_{i\cdot}-\bar
    x_{\cdot\cdot})$ é um estimador não viesado de $\mu_i$, isto é, $E[\bar
    y_{i\cdot(\text{ajustada})}] = \mu_i$.
  \item Em que passo da demonstração o item (b) usaria a hipótese de que a covariável não é
    afetada pelo tratamento? (Não é preciso reescrever a álgebra --- só identifique onde a
    demonstração deixaria de valer se essa hipótese falhasse.)
\end{enumerate}
\end{questao}

\begin{questao}{5} \textbf{(Ciência de dados --- Kruskal-Wallis vs. ANOVA, e o efeito de empates)}

Tempos de carregamento medidos em milissegundos, quando arredondados para o décimo de segundo
mais próximo (uma prática comum em telemetria de produção), frequentemente produzem
\textbf{empates} (valores idênticos) entre observações de grupos diferentes.

\begin{enumerate}[label=(\alph*)]
  \item Explique, em termos gerais, por que o teste de Kruskal-Wallis --- baseado em postos ---
    é uma escolha mais defensável que a ANOVA quando os tempos de carregamento têm distribuição
    fortemente assimétrica à direita (cauda longa de conexões lentas), mesmo com um tamanho de
    amostra razoável.
  \item Quando há muitos empates, o posto atribuído a cada observação empatada é a \emph{média}
    dos postos que ocupariam se não houvesse empate. Isso reduz artificialmente a variância dos
    postos, exigindo uma correção no denominador de $H$ (fator de correção por empates,
    $1 - \sum_j (\tau_j^3-\tau_j)/(N^3-N)$, em que $\tau_j$ é o número de observações no
    $j$-ésimo grupo de empate). O que aconteceria com o valor de $H$ calculado \emph{sem} essa
    correção, na presença de muitos empates: ele tenderia a ficar maior, menor, ou não haveria
    padrão? Justifique intuitivamente.
  \item Kruskal-Wallis testa igualdade de distribuições, não apenas de medianas. Dê um exemplo
    (pode ser qualitativo) de duas distribuições com a mesma mediana, mas que o teste de
    Kruskal-Wallis ainda assim tenderia a distinguir como diferentes.
\end{enumerate}
\end{questao}

\end{document}
